%%%%%%%%%%%%%%%%%%%%%%%%%%%%%%%%%%%%%%%%%%%%
%Bibliothèques de base						%
%%%%%%%%%%%%%%%%%%%%%%%%%%%%%%%%%%%%%%%%%%%%

\documentclass{book}
\pagestyle{headings}
\usepackage[utf8]{inputenc} 					%Encodage des caractères
\usepackage[T1]{fontenc}    					%Encodage de la police

%%%%%%%%%%%%%%%%%%%%%%%%%%%%%%%%%%%%%%%%%%%
%Paquets souvent appelés				   %
%%%%%%%%%%%%%%%%%%%%%%%%%%%%%%%%%%%%%%%%%%%

\usepackage[french]{babel}  					%Règles typographiques françaises
\usepackage{lmodern}							%Choix de la police
\usepackage{amsmath}							%Caractères et environnements mathématiques
\usepackage{amssymb}							%Symboles Mathématiques
\usepackage{amsthm}								%Mise en forme des théorèmes
%\usepackage[left=2cm,right=2cm,top=2cm,bottom=2cm]{geometry}
\usepackage{geometry}
\geometry{a4paper,top=2.5cm,bottom=2cm,inner=3cm,outer=2cm,footskip=1cm}

%%%%%%%%%%%%%%%%%%%%%%%%%%%%%%%%%%%%%%%%%%
%Paquets optionnels	et configuration	  %
%%%%%%%%%%%%%%%%%%%%%%%%%%%%%%%%%%%%%%%%%%
\usepackage[dvipsnames]{xcolor}				%Utiliser des couleurs
\usepackage{ifthen} 							%Faire un peu de tests
\usepackage{tabularx}							%Tableaux avec retour à la ligne
\usepackage{colortbl}							%Couleur dans les tableaux
\usepackage[Sonny]{fncychap}					%Modifier l'en-tête des chapitres
\ChNameVar{\Large}								%Utiliser du serif dans les titres
\ChTitleVar{\Large}								%Utiliser du serif dans les titres
\usepackage{graphicx}

\usepackage{fancyhdr}							%Modifier En-tête et pied de page
\pagestyle{fancy}								%Utiliser un nouveau style de page
\fancyhf{}										%Réinitialiser le style fancy
\fancyhead[LE,RO]{\thepage}						%Style d'en-tête gauche page paire, droit page impaire (extérieur) (ici numéro de page)
\fancyhead[LO]{\nouppercase\rightmark}			%Ajoute le nom du chapitre
\fancyhead[RE]{\nouppercase\leftmark}			%Ajoute le nom de la section
\fancypagestyle{plain}{							%Le style plain est utilisé pour les pages de nouveau chapitre
	\fancyhead{} 								% rien en en-tête
	\renewcommand{\headrulewidth}{0pt}			%filet d'en-tête
	\renewcommand{\footrulewidth}{0pt} 		%filet de pied de page
	\fancyfoot[C]{\thepage}
}

\usepackage{hyperref}							%Utiliser des liens cliquables

\hypersetup{									%Personnaliser hyperref
	%backref=true,								%Permet d'ajouter des liens dans
	%pagebackref=true,							% les bibliographies
	%hyperindex=true,							% Ajoute des liens dans les index.
	colorlinks=true,							% Colorise les liens.(si faux : cadres)
	breaklinks=true,							% Permet le retour à la ligne dans les liens trop longs.
	citecolor=Green,       						% couleur des liens biblios
	filecolor=magenta,      					% couleur des liens vers des fichiers
	urlcolor=blue,           					% couleur des liens externes
	linkcolor= violet,							% Couleur des liens internes.
	%bookmarks=true,								% Créé des signets pour Acrobat.
	bookmarksopen=true,							% Si les signets Acrobat sont créés,
												% les afficher complètement.
	pdftitle={TITRE DE LA THÈSE},		
												% Titre du document.
												% Informations apparaissant dans
	pdfauthor={AUTEUR},					        % dans les informations du document
	pdfsubject={SUJET},			                % sous Acrobat.       
	pdffitwindow=false,     					% Adapter la fenêtre à la page
}



%%%%%%%%%%%%%%%%%%%%%%%%%%%%%%%%%%%%%%%%%%
%Macros personnelles					  %
%%%%%%%%%%%%%%%%%%%%%%%%%%%%%%%%%%%%%%%%%%

\usepackage{lipsum}

%%%%%%%%%%%%%%%%%%%%%%%%%%%%%%%%%%%%%%%%%%
%Informations maketitle					  %
%%%%%%%%%%%%%%%%%%%%%%%%%%%%%%%%%%%%%%%%%%

\title{TITRE DE LA THÈSE}
\author{AUTHER}
\date{DATE SOUTENANCE}
\makeatletter
%%%%%%%%%%%%%%%%%%%%%%%%%%%%%%%%%%%%%%%%%%
%Début du document						  %
%%%%%%%%%%%%%%%%%%%%%%%%%%%%%%%%%%%%%%%%%%

\usepackage{droidsans}
\begin{document}

\begin{titlepage}

	%Définition des données particulières de la page 
	\newcommand{\specialite}{\MakeUppercase{SPÉCIALITÉ}}
	\newcommand{\nomEtab}{l'INSTITUTION} 		%Penser à bien mettre le bon déterminant
	\newcommand{\titreThese}{TITRE THÈSE}
	\newcommand{\nomAuteurice}{Prénom \MakeUppercase{NOM}}
	\newcommand{\dateSoutenance}{DATE}					%Format JJ/MM/AAAA
	
	%%% Jury %%%
	\newcommand{\prefixeUn}{M./Mme}
	\newcommand{\nomJuryUn}{PREMIER MEMBRE}
	\newcommand{\titreJuryUn}{AFFILIATION}
	\newcommand{\roleJuryUn}{RÔLE}
	
	\newcommand{\prefixeDeux}{Mme/M.}
	\newcommand{\nomJuryDeux}{DEUXIÈME MEMBRE}
	\newcommand{\titreJuryDeux}{AFFILIATION PLUS LONGUE POUR FAIRE AU MOINS DEUX LIGNES}
	\newcommand{\roleJuryDeux}{RÔLE}
	
	\newcommand{\prefixeTrois}{}
	\newcommand{\nomJuryTrois}{}
	\newcommand{\titreJuryTrois}{}
	\newcommand{\roleJuryTrois}{}
	
	\newcommand{\prefixeQuatre}{}
	\newcommand{\nomJuryQuatre}{}
	\newcommand{\titreJuryQuatre}{}
	\newcommand{\roleJuryQuatre}{}
	
	\newcommand{\prefixeCinq}{}
	\newcommand{\nomJuryCinq}{}
	\newcommand{\titreJuryCinq}{}
	\newcommand{\roleJuryCinq}{}
	
	\newcommand{\prefixeSix}{}
	\newcommand{\nomJurySix}{}
	\newcommand{\titreJurySix}{}
	\newcommand{\roleJurySix}{}
	
	\newcommand{\prefixeSept}{}
	\newcommand{\nomJurySept}{}
	\newcommand{\titreJurySept}{}
	\newcommand{\roleJurySept}{}
	
	\newcommand{\prefixeHuit}{}
	\newcommand{\nomJuryHuit}{}
	\newcommand{\titreJuryHuit}{}
	\newcommand{\roleJuryHuit}{}
	
	\newcommand{\prefixeNeuf}{}
	\newcommand{\nomJuryNeuf}{}
	\newcommand{\titreJuryNeuf}{}
	\newcommand{\roleJuryNeuf}{}
	
	\newcommand{\nombreMembreJury}{3}
	
	\newcommand{\nomDirecteurUn}{Directeur DE THESE}
	\newcommand{\nomDirecteurDeux}{}
	\newcommand{\nomDirecteurTrois}{}
	
	\newcommand{\nombreDirecteurs}{1}
	
	\newcommand{\nomLabo}{NOM DU LABO (VILLE)}
	
	\droidsans							%Police. Nécessite \usepackage{droidsans}
	\newgeometry{top=1.5cm,bottom=0.2cm,left=1.3cm,right=1.2cm,footskip=1cm}
										%Détermine les dimensions de la page. Nécessite \usepackage{geometry}
										
										%Définition des couleurs. Nécessite \usepackage{xcolor}
	\definecolor{OrangeNU}{HTML}{faa84b}
	\definecolor{VioletNU}{HTML}{650069}
	
	%Pour colorer nécessite \usepackage{colortbl}
	
	\begin{center}
		\includegraphics[scale=1]{logoNU}\\
		\vspace{1cm}
		{\renewcommand{\arraystretch}{2}
		\begin{tabular}{>{\centering\arraybackslash}m{18cm}}
		
		
			\LARGE\textcolor{white}{\textbf{THÈSE}}\cellcolor{VioletNU}\\
			\Large\textbf{Pour obtenir le diplôme de doctorat}\\
			\Large\textbf{Spécialité \specialite}\\
			\Large\textbf{Préparée au sein de \nomEtab}
		\end{tabular}
		\vspace{0.5cm}
		
		}
		{
		\renewcommand{\arraystretch}{2}
		\begin{tabular}{|m{5.8cm}|m{6.4cm}|m{5.7cm}|}
			\multicolumn{3}{m{18cm}}{\centering\LARGE\cellcolor{gray!60}\textbf{\titreThese}}\\
			\multicolumn{3}{c}{\Large\textbf{Présentée et soutenue par}}\\
			\multicolumn{3}{c}{\Large\textbf{\nomAuteurice}}\\
		\end{tabular}
		}
		{
		\renewcommand{\arraystretch}{1.5}
		
		\begin{tabular}{|m{5.8cm}|m{6.4cm}|m{4.9cm}|}
		\arrayrulecolor{OrangeNU}
			\multicolumn{3}{|m{18cm}|}{\centering\cellcolor{OrangeNU}\Large\textbf{Thèse soutenue publiquement le 10/03/2020}}\\
			\hline
			\multicolumn{3}{|m{18cm}|}{\centering\cellcolor{OrangeNU}\Large\textbf{devant le jury composé de}}\\
			\hline
			%%%% Ce passage est particulièrement sale mais j'avais peu de temps %%%%
			\ifnum 0<\nombreMembreJury		
				\prefixeUn\ \MakeUppercase{\nomJuryUn}&\titreJuryUn&\roleJuryUn\\
				\hline
			\fi
			\ifnum 1<\nombreMembreJury		
				\prefixeDeux\ \MakeUppercase{\nomJuryDeux}&\titreJuryDeux&\roleJuryDeux\\
				\hline
			\fi
			\ifnum 2<\nombreMembreJury		
				\prefixeTrois\ \MakeUppercase{\nomJuryTrois}&\titreJuryTrois&\roleJuryTrois\\
				\hline
			\fi
			\ifnum 3<\nombreMembreJury			
				\prefixeQuatre\ \MakeUppercase{\nomJuryQuatre}&\titreJuryQuatre&\roleJuryQuatre\\
				\hline
			\fi
			\ifnum 4<\nombreMembreJury			
				\prefixeCinq\ \MakeUppercase{\nomJuryCinq}&\titreJuryCinq&\roleJuryCinq\\
				\hline
			\fi
			\ifnum 5<\nombreMembreJury			
				\prefixeSix\ \MakeUppercase{\nomJurySix}&\titreJurySix&\roleJurySix\\
				\hline
			\fi
			\ifnum 6<\nombreMembreJury			
				\prefixeSept\ \MakeUppercase{\nomJurySept}&\titreJurySept&\roleJurySept\\
				\hline
			\fi
			\ifnum 7<\nombreMembreJury			
				\prefixeHuit\ \MakeUppercase{\nomJuryHuit}&\titreJuryHuit&\roleJuryHuit\\
				\hline
			\fi
			\ifnum 8<\nombreMembreJury		
				\prefixeNeuf\ \MakeUppercase{\nomJuryNeuf}&\titreJuryNeuf&\roleJuryNeuf\\
				\hline
			\fi
		\end{tabular}		
		 \setlength{\arrayrulewidth}{2mm}{
		 \begin{tabular}{|m{18cm}}
		 \arrayrulecolor{OrangeNU}
		 \ifnum \nombreDirecteurs=1
		 \Large\textbf{Thèse dirigée par \MakeUppercase{\nomDirecteurUn}, \nomLabo}
		 \fi
		 \ifnum \nombreDirecteurs=2
		 \Large\textbf{Thèse dirigée par \MakeUppercase{\nomDirecteurUn} et \MakeUppercase{\nomDirecteurDeux}, \nomLabo}
		 \fi
		 \ifnum \nombreDirecteurs=3
		 \Large\textbf{Thèse dirigée par \MakeUppercase{\nomDirecteurUn}, \MakeUppercase{\nomDirecteurDeux} et \MakeUppercase{\nomDirecteurTrois}, \nomLabo}
		 \fi
		\end{tabular}
		
		}
		}
		
		\vfill
		\begin{tabular}{ccc}
			\includegraphics[width=0.3\textwidth]{logoUni}&\includegraphics[width=0.3\textwidth]{logoED}&\includegraphics[width=0.3\textwidth]{logoLabo}
		\end{tabular}
	\end{center}
\end{titlepage}
\chapter{Titre Premier chapitre}
    \section{Titre Première sections}
    \subsection{Titre Première sous-section}
        \lipsum[1-10]
    \subsection{Titre Deuxième sous-section}
        \lipsum[11-15]
    \section{Titre Deuxième section}
        \lipsum[16-18]
\chapter{Titre Deuxième chapitre}
    \lipsum[19-20]
    \section{Titre Première section}
        \lipsum[21-24]
\end{document}
